\section{Gestion des messages d'erreur}
\subsection{Formatage des erreurs de compilation}
\noindent
Toutes les erreurs détectées avant la génération de code (lexicographie, syntaxe et vérification contextuelle) suivent le formatage requis pour l'évaluation automatique :

\begin{command}
<nom_de_fichier.deca>:<ligne>:<colonne>: <description_de_l_erreur>
\end{command}

\subsection{Erreurs de lexicographie (Lexer)}
\noindent
L'analyseur lexical se charge de reconnaître les jetons (\texttt{tokens}). Les erreurs surviennent lorsqu'une séquence ne correspond à aucune règle définie. Ces erreurs sont capturées par le \texttt{DecacErrorListener} ou de rattrapées avant dans le \texttt{DecaLexer} lui-même (pour \texttt{INVALID\_IDENT} et \texttt{UNTERMINATED\_COMMENT}) pour fournir des messages plus clairs:


\begin{itemize}[label=$\circ$]
    \item \texttt{INVALID\_IDENT} : Levé par un identificateur qui ne respecte pas les règles lexicographiques (ex : commence par un chiffre).
    \item \texttt{UNTERMINATED\_COMMENT} : Levé lorsqu'un commentaire multi-ligne (\texttt{/*}) est ouvert mais jamais refermé avant la fin du fichier.
    \item \texttt{token recognition error} : Levé lors de l'utilisation de caractères n'appartenant pas au langage Deca (ex : \texttt{\#}, \texttt{@}) ou d'une chaîne de caractères mal formée (ex : guillemet manquant).
\end{itemize}

\subsection{Erreurs de syntaxe hors-contexte (Parser)}
\noindent
L'analyseur syntaxique vérifie si l'ordre des jetons respecte la grammaire du langage. Sous chaque explication d'erreur est fourni un exemple de fichier test qui lève l'erreur. Ces erreurs sont capturées par le \texttt{DecacErrorListener} :

\begin{itemize}[label=$\circ$]
    \item \texttt{mismatched input} : Un jeton a été trouvé là où un autre était attendu (ex : condition de \texttt{if} vide).
    \newline {\footnotesize Exemple de fichier test : \texttt{.../syntax/invalid/parser/empty-if.deca}}
    
    \item \texttt{no viable alternative} : Structure de code incomplète ou ambiguë (ex : point-virgule manquant).
    \newline {\footnotesize Exemple de fichier test : \texttt{.../syntax/invalid/parser/missing-semi.deca}}
    
    \item \texttt{missing <token>} : Absence d'un symbole obligatoire (ex : parenthèse fermante oubliée).
    \newline {\footnotesize Exemple de fichier test : \texttt{.../syntax/invalid/parser/missing-token.deca}}
    
    \item \texttt{extraneous input} : Symbole invalide pour la règle en cours (ex : opérateur doublé).
    \newline {\footnotesize Exemple de fichier test : \texttt{.../syntax/invalid/parser/duplicate-op.deca}}
\end{itemize}


\subsection{Erreurs de syntaxe contextuelle}
\noindent
Ces erreurs lors de la décoration de l'AST après qu'il soit construit. Elles concernent la cohérence des types, des variables et des déclarations.

\begin{table}[H]
\centering
\begin{tabularx}{\textwidth}{|l|X|}
\hline
\rowcolor{gray!10} Catégorie & Messages d'erreur et causes possibles \\
\hline
Identificateurs & \texttt{Identifier <name> not found} : variable ou méthode non définie. \newline \texttt{Type <name> is already defined} : double déclaration d'un type. \\
\hline
Classes & \texttt{Superclass <name> of class <C> is not defined} : héritage d'une classe inexistante. \newline \texttt{New can only be used with class types} : instanciation sur un type primitif. \\
\hline
Méthodes/Champs & \texttt{Field name collides with an inherited method} : conflit lors de l'héritage. \newline \texttt{Visibility error for protected field} : accès illégal à un champ protégé. \\
\hline
Types/Opérations & \texttt{Operand of UnaryMinus must be of numeric type} : \texttt{-} appliqué sur un booléen. \newline \texttt{Condition must be a boolean. Found: <type>} : type invalide dans un \texttt{if} ou \texttt{while}. \newline \texttt{Cannot cast from <T1> to <T2>} : cast illégal entre types incompatibles. \\
\hline
Arguments & \texttt{Expected <n> args, got <m>} : signature de méthode non respectée. \\
\hline
\end{tabularx}
\end{table}

\subsection{Erreurs d'exécution (Runtime)}
\noindent
Même si le programme est valide contextuellement, il peut encore échouer lors de l'exécution sur la machine virtuelle IMA (règles sémantiques). Ces erreurs sont gérées par notre \texttt{RuntimeErrorManager}. Chaque erreur d'exécution est liée à un label spécifique en assembleur. Lorsqu'une condition d'erreur est détectée (avec l'instruction \texttt{BOV} par exemple), le programme passe à l'instruction du label ce qui affiche le message et stoppe l'exécution (\texttt{ERROR}).

\begin{table}[H]
\centering
\begin{tabularx}{\textwidth}{|l|X|}
\hline
\rowcolor{gray!10} Message d'erreur & Cause(s) possible(s) de l'erreur \\
\hline
\texttt{Division by zero} & Tentative de division par un entier ou flottant nul. \\
\hline
\texttt{Stack overflow} & Débordement de la pile (appels récursifs trop profonds ou registres). \\
\hline
\texttt{Input/Output error} & Erreur lors de la lecture d'un entier ou d'un flottant \\
\hline
\texttt{Overflow arith.} & Débordement lors d'une opération sur les nombres flottants. \\
\hline
\texttt{Null dereference} & Accès à un champ ou une méthode sur un objet non instancié. \\
\hline
\texttt{Heap overflow} & Mémoire globale saturée lors d'une allocation par \texttt{new}. \\
\hline
\texttt{Impossible cast} & Échec d'un transtypage dynamique entre deux classes incompatibles. \\
\hline
\end{tabularx}
\end{table}